\section{Domain Task 8: Phân tích mùa vụ và cơ hội đặt phòng}

Câu hỏi: Tỷ lệ lấp đầy thay đổi như thế nào theo mùa? Chính sách minimum
nights ảnh hưởng thế nào?

\subsection{Biểu đồ 1 -- Line Chart: Tỷ lệ lấp đầy theo tháng}

\subsubsection*{A. Thiết kế Idiom}

\begin{longtable}{@{} p{2.8cm} p{12cm} @{}}
\toprule
\textbf{Đặc điểm} & \textbf{Chi tiết} \\
\midrule
\endhead

\bottomrule
\endfoot

\textbf{Idiom} & Line Chart \\
\addlinespace[0.2cm]

\textbf{What} & Month: Ordinal (tháng 1--12, cyclic -- seasonality lặp lại theo năm) \newline
occupancy\_rate\_pct (derived): Quantitative (AVG(is\_booked) $\times$ 100) \\
\addlinespace[0.2cm]

\textbf{How} & \textbf{Encode:}
\begin{itemize}[nosep, leftmargin=*, topsep=0.1cm]
    \item Mark: Line + Point
    \item Channel:
    \begin{itemize}[nosep, leftmargin=0.5cm]
        \item Pos X: Month (Ordinal, Discrete -- tiến trình thời gian tháng 1--12)
        \item Pos Y: Occupancy Rate (\%) (Quantitative)
        \item Reference Line: Average (ngang -- mức trung bình tổng thể $\sim$30\%)
    \end{itemize}
\end{itemize}

\vspace{0.2cm}
\textbf{Manipulate:}
\begin{itemize}[nosep, leftmargin=*, topsep=0.1cm]
    \item Hover để xem tháng, occupancy rate (\%) cụ thể
\end{itemize}

\vspace{0.2cm}
\textbf{Facet:}
\begin{itemize}[nosep, leftmargin=*, topsep=0.1cm]
    \item Superimpose: Reference Line (average occupancy ngang)
\end{itemize}

\vspace{0.2cm}
\textbf{Reduce:}
\begin{itemize}[nosep, leftmargin=*, topsep=0.1cm]
    \item Filter: Year = 2026 (dữ liệu calendar 2025 chỉ có từ tháng 11, không đủ để phân tích seasonality)
\end{itemize} \\
\addlinespace[0.2cm]

\textbf{Why} & discover $\rightarrow$ browse \newline
Line chart (không phải bar) vì Month là Ordinal có thứ tự -- đường kết nối nhấn mạnh tính liên tục và xu hướng thời gian, không phải giá trị tại từng điểm rời rạc. \\
\addlinespace[0.2cm]

\textbf{Scale} &
\begin{itemize}[nosep, leftmargin=*, topsep=0pt]
    \item Key: 12 (tháng); Items: dữ liệu 2026 (tháng 1--11)
\end{itemize} \\
\end{longtable}

\begin{figure}[H]
    \centering
    \safeincludegraphics[width=0.9\linewidth]{charts/domain_task8_chart1.png}
    \caption{Hình 8.1. Tỷ lệ lấp đầy (occupancy rate) theo tháng tại NYC (2026)}
    \label{fig:domain-task8-chart1}
\end{figure}

\subsubsection*{B. Đánh giá biểu đồ}

\textbf{1. Nguyên lý biểu đạt (Expressiveness):}

\begin{itemize}
    \item Thuộc tính: Tháng / Month (O), Tỷ lệ lấp đầy / Occupancy Rate Pct (Q -- AVG aggregate), Đường trung bình / Average Occupancy Rate (Q -- Reference Line), Năm / Year (O -- filter bối cảnh, Year = 2026).
    \item Channel:
    \begin{itemize}
        \item PosX (Month -- Ordinal): trục thời gian có thứ tự $\rightarrow$ đúng. Không dùng Categorical vì mất thứ tự thời gian.
        \item PosY (Occupancy Rate -- Q): thể hiện mức độ lấp đầy $\rightarrow$ đúng.
        \item Reference Line Average: cung cấp mốc tham chiếu tổng thể $\rightarrow$ tăng thêm thông tin phân tích.
    \end{itemize}
    \item Đánh giá mức độ quan trọng:
    \begin{itemize}
        \item Ưu tiên 1 -- Vị trí dọc (PosY --- Position on common scale): PosY là kênh quan trọng nhất trong line chart. Kênh này mã hóa Occupancy Rate (\%) (Quantitative) trên thang đo chung, phản ánh trực tiếp mức độ lấp đầy của thị trường tại từng thời điểm. Theo Mackinlay, Position on a common scale là kênh hiệu quả nhất cho dữ liệu định lượng, cho phép người xem so sánh mức occupancy giữa các tháng với độ chính xác cao và nhận diện biên độ dao động mùa vụ (từ $\sim$20\% đến $\sim$40\%) một cách tức thì. Kênh này là nền tảng để thực hiện task discover (phát hiện xu hướng) và browse (duyệt qua từng tháng).
        \item Ưu tiên 2 -- Vị trí ngang (PosX --- trục thời gian Ordinal): PosX mã hóa Month (Ordinal --- tháng 1 đến 12) theo trục thời gian. Đây là kênh quan trọng thứ hai vì nó thiết lập chiều thời gian của biểu đồ --- nền tảng của phân tích seasonality. Theo Mackinlay, Position phù hợp cho dữ liệu thứ tự (Ordinal), và việc sử dụng trục ngang liên tục phản ánh đúng bản chất thứ tự tháng (T1 $\rightarrow$ T12), không bị mất thông tin thứ tự như khi dùng Categorical. Các điểm được kết nối bằng đường thẳng giúp người xem nhận diện xu hướng liên tục theo thời gian.
        \item Ưu tiên 3 -- Reference Line (Average): Reference line ngang tại mức trung bình ($\sim$30\%) là channel bổ trợ giúp người xem định vị nhanh các tháng vượt/dưới ngưỡng bình thường mà không cần đọc chính xác giá trị trục Y. Đây là annotation tĩnh (không encode data mới) nhưng tăng đáng kể khả năng thực hiện action Locate Extremes --- xác định tháng cao điểm và thấp điểm một cách tức thì. Theo phân loại Munzner, đây là dạng derived encoding bổ trợ cho channel chính thay vì thay thế nó.
    \end{itemize}
    \item Kết luận: Line chart là idiom tối ưu cho dữ liệu chuỗi thời gian, đặc biệt khi cần phát hiện trend và seasonality.
\end{itemize}

\textbf{2. Nguyên lý hiệu quả (Effectiveness):}

\begin{itemize}
    \item Độ chính xác (Accuracy): PosY (position on common scale) --- rất chính xác. Người xem dễ dàng so sánh mức occupancy giữa các tháng.
    \item Discriminability: Reference line Average tạo mốc tham chiếu để xác định tháng nào trên/dưới trung bình.
    \item Separability: PosX và PosY tách biệt hoàn toàn.
\end{itemize}

\subsubsection*{C. Phân tích biểu đồ (Insight)}

\begin{itemize}
    \item Tháng 1--4 là low season rõ rệt (occupancy $\sim$20--21\%) --- cơ hội tốt cho khách muốn giá thấp và dễ tìm phòng.
    \item Từ tháng 5 trở đi, occupancy tăng dần và đạt đỉnh vào tháng 11 ($\sim$40\%) --- phản ánh nhu cầu du lịch mùa thu và dịp lễ Thanksgiving.
    \item Dữ liệu sử dụng: calendar 2026 (tháng 1--11); dữ liệu 2025 chỉ có từ tháng 11 nên không đủ để so sánh liên năm.
    \item Khuyến nghị cho du khách: Đặt phòng tháng 1--3 để có giá thấp nhất; tháng 11 cần đặt sớm do cầu rất cao.
    \item Khuyến nghị cho host: Tháng 1--4 nên linh hoạt hóa chính sách (giảm minimum nights, giảm giá nhẹ); tháng 11--12 có thể tăng giá.
\end{itemize}
